\chapter{Vecteurs propres, valeurs propres et diagonalisation}
\label{app:vecteurs-propres-diagonalisation}

Soit $V$ un espace vectoriel de dimension finie $N$ sur un corps $\mathbb{K}$, et soit
\[
f : V \rightarrow V
\]
un endomorphisme, c'est-à-dire une application linéaire de $V$ dans lui-même.

Un vecteur non nul $v \in V$ est appelé vecteur propre de $f$ s'il existe un scalaire
$\lambda \in \mathbb{K}$ tel que :
\[
f(v) = \lambda v.
\]
Le scalaire $\lambda$ est alors appelé valeur propre associée au vecteur propre $v$.

Cette relation signifie que l'application $f$ ne change pas la direction du vecteur $v$.
Elle se contente de le multiplier par un facteur scalaire $\lambda$.

Considérons maintenant une base
\[
\mathcal{B} = (e_1,\dots,e_N)
\]
de $V$. D'après le rappel de l'annexe~\ref{app:representation-lineaire-matrice},
l'endomorphisme $f$ peut être représenté par une matrice dans cette base.

La question naturelle est alors la suivante : peut-on choisir une base dans laquelle la matrice
de $f$ devient particulièrement simple ?

Le cas le plus favorable est celui où la base $\mathcal{B}$ est composée uniquement de
vecteurs propres de $f$. Supposons donc que, pour tout $i \in \{1,\dots,N\}$, on ait :
\[
f(e_i) = \lambda_i e_i.
\]

Alors, dans la base $\mathcal{B}$, les coordonnées de $f(e_i)$ sont nulles partout sauf à la
$i$-ème position :
\[
[f(e_i)]_{\mathcal{B}}
=
\begin{pmatrix}
0 \\
\vdots \\
0 \\
\lambda_i \\
0 \\
\vdots \\
0
\end{pmatrix}.
\]

La matrice représentant $f$ dans cette base a donc pour colonnes ces vecteurs de coordonnées.
Elle est diagonale :
\[
A_{\mathcal{B}}
=
\begin{pmatrix}
\lambda_1 & 0         & \cdots & 0 \\
0         & \lambda_2 & \cdots & 0 \\
\vdots    & \vdots    & \ddots & \vdots \\
0         & 0         & \cdots & \lambda_N
\end{pmatrix}.
\]

Autrement dit, si l'on dispose d'une base de vecteurs propres, l'action de $f$ devient très simple :
chaque coordonnée est simplement multipliée par la valeur propre correspondante.

Réciproquement, si la matrice de $f$ est diagonale dans une base donnée, alors les vecteurs de cette
base sont des vecteurs propres de $f$.

Cette observation explique l'intérêt de chercher une base de vecteurs propres : dans une telle base,
l'application linéaire ne mélange plus les directions. Elle agit indépendamment sur chaque direction
de la base.

Dans le cadre de ce rapport, cette idée est appliquée à la convolution circulaire. La matrice associée
à une convolution circulaire admet les vecteurs de Fourier comme vecteurs propres. Ainsi, lorsque
l'on exprime le signal dans la base de Fourier, l'opérateur de convolution devient diagonal :
\[
C_h = F_N \Lambda_h F_N^{-1},
\]
où $F_N$ contient les vecteurs de Fourier et où $\Lambda_h$ est la matrice diagonale des valeurs
propres associées.